\section{Literature Review}
Raskar et al.\ applied temporal coding using a fluttered shutter to encode motion within the exposure and deblur images\cite{raskar}. We draw inspiration for the use of temporal coding, implementing the method using illumination and modeling axial scaling of the image. Our binary pattern is not optimized and is quite short. Their lateral-motion results are motivating to conduct the experiment, but do not constitute a demonstration that coding will benefit this apparatus.

We apply a total variation (TV) regularization to keep noise down, while still preserving edges\cite{rof}. The solver is an accelerated proximal-gradient algorithm based on FISTA\cite{fista}. It involves an approximate TV proximal step, for which we check numerical stability. Endpoint registration is based on scale-invariant features\cite{lowe}. Pixel error is measured with peak signal-to-noise ratio (PSNR), and the quality of local structure is measured by structural similarity (SSIM)\cite{ssim}. The inclusion of the blurred input image as a baseline provides a check of whether the resulting reconstruction helps to make a better photo in addition to the assumed motion model.
